% !TEX root = ../main.tex
%% Corrigés — Chapitre 16 — Dénombrement
%% Résultats vérifiés par verifications/chapitre_16.py.

\corrigesChapitre{16}{Dénombrement}
\corrigesSouples

\begin{blocLivre}{cor-16-1}
\begin{corrige}{1}{Premiers décomptes}
\demarche{Des choix successifs indépendants : on multiplie.}
\begin{questions}
  \item $3\times4\times2=\res{24}$ menus.
  \item $10^{4}=\res{10\,000}$ codes (les répétitions sont permises).
  \item $26^{2}\times10^{3}\times26^{2}=\res{456\,976\,000}$ plaques.
\end{questions}
\end{corrige}
\end{blocLivre}

\begin{blocLivre}{cor-16-2}
\begin{corrige}{2}{Unions et complémentaires}
\demarche{Additionner, puis retrancher ce qui a été compté deux fois.}
\begin{questions}
  \item $|M\cup P|=18+15-8=\res{25}$ élèves suivent au moins une option ;
    $30-25=\res{5}$ n'en suivent aucune.
  \item $A\cup B$ est l'union \emph{disjointe} de $A$ et de
    $B\setminus A$, et $|B\setminus A|=|B|-|A\cap B|$ : en sommant,
    $|A\cup B|=|A|+|B|-|A\cap B|$ — l'intersection, comptée dans $|A|$ et
    dans $|B|$, est retranchée une fois.
\end{questions}
\end{corrige}
\end{blocLivre}

\begin{blocLivre}{cor-16-3}
\begin{corrige}{3}{Listes, avec ou sans répétition}
\demarche{Répétition permise : une puissance. Interdite : une factorielle
descendante.}
\begin{questions}
  \item $26^{3}=\res{17\,576}$ mots ; sans répétition :
    $26\times25\times24=\res{15\,600}$.
  \item $5^{3}=\res{125}$ applications ; injectives :
    $5\times4\times3=\res{60}$ — à chaque étape, une image de moins est
    disponible.
\end{questions}
\end{corrige}
\end{blocLivre}

\begin{blocLivre}{cor-16-4}
\begin{corrige}{4}{Permutations et anagrammes}
\demarche{Des lettres répétées : diviser par leurs permutations internes.}
\begin{questions}
  \item $4!=\res{24}$ anagrammes de MPSI ; $5!=\res{120}$ rangements.
  \item En distinguant fictivement $\mathrm{T}_1$ et $\mathrm{T}_2$ :
    $7!$ listes ; chaque anagramme réelle est obtenue $2!$ fois (échange
    des deux T) : $\dfrac{7!}{2!}=\res{2520}$.
\end{questions}
\end{corrige}
\end{blocLivre}

\begin{blocLivre}{cor-16-5}
\begin{corrige}{5}{Combinaisons : premiers calculs}
\demarche{Une main, un comité : des parties, pas des listes.}
\begin{questions}
  \item $\dbinom52=\res{10}$, $\dbinom{10}3=\res{120}$,
    $\dbinom{12}4=\res{495}$.
  \item $\dbinom{32}5=\res{201\,376}$ mains.
  \item Une main est un \emph{ensemble} de cartes : les
    $\frac{32!}{27!}$ tirages ordonnés comptent chaque main $5!$ fois,
    d'où la division par $5!$ dans la formule.
\end{questions}
\end{corrige}
\end{blocLivre}

\begin{blocLivre}{cor-16-6}
\begin{corrige}{6}{Le triangle de Pascal}
\demarche{Chaque case est la somme des deux cases au-dessus — et chaque
ligne se lit comme des comptes de parties.}
\begin{questions}
  \item Lignes successives :
    $1$ ; $1\,1$ ; $1\,2\,1$ ; $1\,3\,3\,1$ ; $1\,4\,6\,4\,1$ ;
    $1\,5\,10\,10\,5\,1$ ; $1\,6\,15\,20\,15\,6\,1$.
  \item $\dbinom63=\res{20}$ ; la ligne $6$ est un palindrome
    ($\binom np=\binom n{n-p}$), et sa somme vaut $64=2^{6}$.
  \item Fixons $a\in E$, $|E|=n$ : les parties à $p$ éléments contenant
    $a$ sont au nombre de $\binom{n-1}{p-1}$, celles ne le contenant pas
    $\binom{n-1}p$ : la partition donne
    $\binom np=\binom{n-1}{p-1}+\binom{n-1}{p}$.
\end{questions}
\end{corrige}
\end{blocLivre}

\begin{blocLivre}{cor-16-7}
\begin{corrige}{7}{Le binôme en action}
\demarche{Des valeurs particulières de a et b font jaillir les identités.}
\begin{questions}
  \item $(a+b)^{5}=a^{5}+5a^{4}b+10a^{3}b^{2}+10a^{2}b^{3}+5ab^{4}+b^{5}$ ;
    le coefficient de $x^{3}$ dans $(1+x)^{7}$ est
    $\dbinom73=\res{35}$.
  \item $a=b=1$ : $\displaystyle\sum_k\binom nk=2^{n}$ ;
    $a=-1$, $b=1$ :
    $\displaystyle\sum_k(-1)^{k}\binom nk=0$ pour $n\geqslant1$.
\end{questions}
\end{corrige}
\end{blocLivre}

\begin{blocLivre}{cor-16-8}
\begin{corrige}{8}{Compter des parties}
\demarche{Une partie = une suite de n décisions binaires.}
\begin{questions}
  \item L'application $A\mapsto\ind{A}$ est une bijection de
    $\parties{E}$ sur les applications de $E$ dans $\ens{0,1}$ :
    $|\parties{E}|=2^{n}$.
  \item Une partie contenant $a$ est déterminée par sa trace sur
    $E\setminus\ens{a}$ : $\res{2^{n-1}}$.
  \item La somme alternée $\sum(-1)^{k}\binom nk=0$ dit que
    $\sum_{k\text{ pair}}\binom nk=\sum_{k\text{ impair}}\binom nk$ :
    chaque camp compte $\res{2^{n-1}}$ parties.
\end{questions}
\end{corrige}
\end{blocLivre}

\begin{blocLivre}{cor-16-9}
\begin{corrige}{9}{Le principe des tiroirs}
\demarche{Plus d'objets que de tiroirs : un tiroir déborde. Tout l'art est
de choisir les tiroirs.}
\begin{questions}
  \item $13$ personnes, $12$ mois : deux personnes partagent un mois.
  \item Tiroirs : les $n$ paires $\ens{1,2},\ens{3,4},\dots,
    \ens{2n-1,2n}$. Avec $n+1$ entiers, deux tombent dans la même paire :
    ils sont consécutifs, et $d\mid a$ et $d\mid a+1$ forcent $d\mid1$ :
    premiers entre eux (chapitre 7).
  \item Découper le carré en $4$ carrés de côté $\frac12$, de diagonale
    $\frac{\sqrt2}{2}$ : avec $5$ points, deux sont dans le même petit
    carré, donc à distance au plus $\res{\dfrac{\sqrt2}{2}}$.
\end{questions}
\end{corrige}
\end{blocLivre}

\begin{blocLivre}{cor-16-10}
\begin{corrige}{10}{Comités et délégations}
\demarche{« Exactement » : un produit. « Au moins » : le complémentaire.
Deux comptes d'un même objet : une identité.}
\begin{questions}
  \item Choisir $2$ femmes et $2$ hommes :
    $\dbinom72\dbinom52=21\times10=\res{210}$.
  \item Complémentaire (comités sans femme) :
    $\dbinom{12}4-\dbinom54=495-5=\res{490}$.
  \item Comité puis président : $\dbinom{12}4\times4$ ; président puis
    reste du comité : $12\times\dbinom{11}3$. Les deux valent
    $\res{1980}$ : c'est l'identité $k\binom nk=n\binom{n-1}{k-1}$ pour
    $n=12$, $k=4$ (exercice 13).
\end{questions}
\end{corrige}
\end{blocLivre}

\begin{blocLivre}{cor-16-11}
\begin{corrige}{11}{Mains de cartes}
\demarche{Partitionner la main : les as d'un côté, le reste de l'autre.}
\begin{questions}
  \item $2$ as parmi $4$, puis $3$ cartes parmi les $28$ autres :
    $\dbinom42\dbinom{28}3=6\times3276=\res{19\,656}$.
  \item Complémentaire (aucun as, soit $5$ cartes parmi $28$) :
    $\dbinom{32}5-\dbinom{28}5=201\,376-98\,280=\res{103\,096}$.
\end{questions}
\end{corrige}
\end{blocLivre}

\begin{blocLivre}{cor-16-12}
\begin{corrige}{12}{Chemins sur une grille}
\demarche{Un chemin est un mot : il suffit de placer les pas « droite ».}
\begin{questions}
  \item Un chemin est une suite de $p+q$ pas dont $p$ « droite » : le
    choix de leurs positions donne $\res{\dbinom{p+q}{p}}$ chemins.
  \item $\dbinom74=\res{35}$ ; par $(1,2)$ : $\dbinom31$ chemins jusqu'à
    $(1,2)$ puis $\dbinom43$ jusqu'à $(4,3)$, soit $3\times4=\res{12}$.
  \item Le dernier pas arrive de $(p-1,q)$ ou de $(p,q-1)$ :
    $\binom{p+q}{p}=\binom{p+q-1}{p-1}+\binom{p+q-1}{p}$ — la formule de
    Pascal, dessinée sur la grille.
\end{questions}
\medskip
\begin{center}
\begin{tikzpicture}[scale=0.72,line cap=round]
  \draw[otMuted!60,line width=0.4pt] (0,0) grid (4,3);
  \draw[otGreenDark,line width=1.6pt,line join=round]
    (0,0) -- (1,0) -- (1,2) -- (3,2) -- (3,3) -- (4,3);
  \fill[otNavy] (1,2) circle (2.2pt);
  \fill[otInk] (0,0) circle (2pt);
  \fill[otInk] (4,3) circle (2pt);
  \node[below left,font=\scriptsize,text=otInk] at (0,0) {$(0,0)$};
  \node[above right,font=\scriptsize,text=otInk] at (4,3) {$(4,3)$};
  \node[above left,font=\scriptsize,text=otNavy] at (1,2) {$(1,2)$};
\end{tikzpicture}\\[0.5mm]
\legende{Un des $12$ chemins de $(0,0)$ à $(4,3)$ passant par $(1,2)$ :
choisir les positions des pas « droite » suffit.}
\end{center}
\end{corrige}
\end{blocLivre}

\begin{blocLivre}{cor-16-13}
\begin{corrige}{13}{Le comité et son président}
\demarche{Compter les comités présidés de deux façons, puis sommer sur la
taille.}
\begin{questions}
  \item Comité de $k$ puis président dedans : $\binom nk\,k$ ; président
    d'abord puis $k-1$ collègues parmi $n-1$ : $n\binom{n-1}{k-1}$.
    D'où $\res{k\dbinom nk=n\dbinom{n-1}{k-1}}$.
  \item En sommant sur $k$ : un « comité présidé de taille quelconque »
    se construit en choisissant le président ($n$ façons) puis une partie
    quelconque du reste ($2^{n-1}$) :
    $\sum_k k\binom nk=\res{n\,2^{n-1}}$. Avec président et trésorier
    distincts : $n(n-1)$ choix puis $2^{n-2}$ :
    $\sum_k k(k-1)\binom nk=\res{n(n-1)\,2^{n-2}}$.
\end{questions}
\end{corrige}
\end{blocLivre}

\begin{blocLivre}{cor-16-14}
\begin{corrige}{14}{La formule de Vandermonde}
\demarche{Partitionner les choix selon la part prise dans chaque groupe.}
\begin{questions}
  \item Pour choisir $n$ personnes parmi $p$ femmes et $q$ hommes, on
    classe selon le nombre $k$ de femmes retenues :
    $\binom pk$ choix de femmes et $\binom q{n-k}$ d'hommes, d'où
    $\displaystyle\sum_{k}\binom pk\binom q{n-k}=\binom{p+q}{n}$.
  \item Avec $p=q=n$ et la symétrie $\binom n{n-k}=\binom nk$ :
    \[
    \res{\sum_{k=0}^{n}\binom nk^{2}=\binom{2n}{n}}
    \]
    — le binôme central, déjà croisé aux chapitres 11 et 15 avec Wallis
    et Stirling.
\end{questions}
\end{corrige}
\end{blocLivre}

\begin{blocLivre}{cor-16-15}
\begin{corrige}{15}{Surjections sur deux ou trois éléments}
\demarche{Compter ce qui manque : les applications qui évitent une valeur.}
\begin{questions}
  \item Parmi les $2^{n}$ applications, seules les deux constantes ne
    sont pas surjectives : $\res{2^{n}-2}$.
  \item Soit $A_i$ l'ensemble des applications évitant la valeur $i$ :
    $|A_i|=2^{n}$, $|A_i\cap A_j|=1$ (constantes), triple intersection
    vide. Deux applications de la formule à deux ensembles donnent
    $|A_1\cup A_2\cup A_3|=3\cdot2^{n}-3$, d'où
    $\res{3^{n}-3\cdot2^{n}+3}$ surjections — le crible en action
    (exercice 20).
\end{questions}
\end{corrige}
\end{blocLivre}

\begin{blocLivre}{cor-16-16}
\begin{corrige}{16}{Étoiles et barres}
\demarche{Coder chaque solution par un dessin : des étoiles séparées par
des barres.}
\begin{questions}
  \item Alignons $n$ étoiles et $p-1$ barres : les barres découpent les
    étoiles en $p$ paquets de tailles $x_1,\dots,x_p$ (éventuellement
    vides). Ce codage est une bijection : une solution correspond
    exactement à un choix des positions des $p-1$ barres parmi les
    $n+p-1$ symboles :
    $\res{\dbinom{n+p-1}{p-1}}$ solutions.
  \item $x+y+z=10$ : $\dbinom{12}2=\res{66}$.
\end{questions}
\end{corrige}
\end{blocLivre}

\begin{blocLivre}{cor-16-17}
\begin{corrige}{17}{Applications monotones}
\demarche{Strictement croissante = une partie ; croissante = un décalage
pour écarter les égalités.}
\begin{questions}
  \item Une application strictement croissante est entièrement déterminée
    par son image, partie à $p$ éléments de $\ens{1,\dots,n}$ — et toute
    partie, rangée en ordre croissant, en fournit une :
    $\res{\dbinom np}$.
  \item Poser $g(i)=f(i)+i-1$ : $f$ est croissante si et seulement si
    $g$ est \emph{strictement} croissante de $\ens{1,\dots,p}$ dans
    $\ens{1,\dots,n+p-1}$ : d'après a),
    $\res{\dbinom{n+p-1}{p}}$ applications — le compte « avec répétition,
    sans ordre » qui complète le tableau du cours.
\end{questions}
\end{corrige}
\end{blocLivre}

\begin{blocLivre}{cor-16-18}
\begin{corrige}{18}{Le binôme rencontre l'analyse}
\demarche{Intégrer ou dériver l'identité du binôme terme à terme.}
\begin{questions}
  \item $\displaystyle\int_{0}^{1}(1+x)^{n}\mathrm{d}x
    =\frac{2^{n+1}-1}{n+1}$, et terme à terme
    $\displaystyle\int_{0}^{1}\sum_k\binom nk x^{k}
    =\sum_k\binom nk\frac{1}{k+1}$ :
    \[
    \res{\sum_{k=0}^{n}\frac{1}{k+1}\binom nk=\frac{2^{n+1}-1}{n+1}}.
    \]
  \item En dérivant : $n(1+x)^{n-1}=\sum_k k\binom nk x^{k-1}$, et
    $x=1$ redonne $\sum_k k\binom nk=n\,2^{n-1}$ — l'analyse confirme le
    double comptage de l'exercice 13.
\end{questions}
\end{corrige}
\end{blocLivre}

\begin{blocLivre}{cor-16-19}
\begin{corrige}{19}{Les dérangements}
\demarche{Classer selon le sort du dernier élément, puis télescoper.}
\begin{questions}
  \item Soit $\sigma$ un dérangement de $\ens{1,\dots,n}$ et
    $j=\sigma(n)\neq n$ ($n-1$ choix). Si $\sigma(j)=n$, le reste est un
    dérangement des $n-2$ autres éléments : $D_{n-2}$ cas. Sinon, en
    court-circuitant $n$ (envoyer l'antécédent de $n$ directement sur
    $j$), on obtient un dérangement de $n-1$ éléments : $D_{n-1}$ cas.
    D'où $\res{D_n=(n-1)\bigl(D_{n-1}+D_{n-2}\bigr)}$.
  \item En posant $E_n=D_n-nD_{n-1}$, la relation donne
    $E_n=-E_{n-1}$ avec $E_1=-1$ : $E_n=(-1)^{n}$, soit
    $D_n=nD_{n-1}+(-1)^{n}$. En divisant par $n!$ :
    $\dfrac{D_n}{n!}=\dfrac{D_{n-1}}{(n-1)!}+\dfrac{(-1)^{n}}{n!}$, et la
    somme télescopique donne
    $\res{D_n=n!\,\displaystyle\sum_{k=0}^{n}\frac{(-1)^{k}}{k!}}$.
  \item La somme partielle tend vers
    $\mathrm{e}^{-1}$ (chapitres 5 et 15) :
    $\dfrac{D_n}{n!}\to\res{\dfrac{1}{\mathrm{e}}}\approx0{,}368$ — plus
    d'un tiers des permutations ne fixent rien.
\end{questions}
\end{corrige}
\end{blocLivre}

\begin{blocLivre}{cor-16-20}
\begin{corrige}{20}{Le crible pour trois ensembles}
\demarche{Appliquer deux fois la formule à deux ensembles, en distribuant
l'intersection.}
\begin{questions}
  \item $|A\cup B\cup C|=|A\cup B|+|C|-|(A\cup B)\cap C|$, et
    $(A\cup B)\cap C=(A\cap C)\cup(B\cap C)$ se développe à son tour :
    en recollant,
    \[
    |A\cup B\cup C|=|A|+|B|+|C|-|A\cap B|-|A\cap C|-|B\cap C|
    +|A\cap B\cap C| .
    \]
  \item Multiples de $2,3,5$ jusqu'à $100$ : $50,33,20$ ; de
    $6,10,15$ : $16,10,6$ ; de $30$ : $3$. L'union compte
    $50+33+20-16-10-6+3=74$ : il reste $\res{26}$ entiers divisibles par
    aucun des trois — les « survivants » du crible d'Ératosthène
    (chapitre 7).
\end{questions}
\end{corrige}
\end{blocLivre}

\begin{blocLivre}{cor-16-21}
\begin{corrige}{21}{Chemins sous la diagonale : les nombres de Catalan}
\demarche{Réfléchir le début du chemin fautif : les fautifs deviennent des
chemins ordinaires vers un autre coin.}
\begin{questions}
  \item À la main : $C_1=1$, $C_2=2$, $C_3=\res{5}$.
  \item Un chemin fautif touche la droite $y=x+1$ ; en réfléchissant par
    rapport à elle la portion \emph{avant} le premier contact, on obtient
    un chemin de $(-1,1)$ — autrement dit un chemin ordinaire compté par
    $\dbinom{2n}{n-1}$ — et cette correspondance est une bijection. D'où
    \[
    C_n=\binom{2n}{n}-\binom{2n}{n-1}
    =\res{\frac{1}{n+1}\binom{2n}{n}}
    \]
    (le calcul de la différence est immédiat avec les factorielles).
  \item $1,\ 2,\ 5,\ 14,\ 42,\dots$ — ces nombres compteront aussi les
    parenthésages et bien d'autres structures.
\end{questions}
\medskip
\begin{center}
\begin{tikzpicture}[scale=0.82,line cap=round]
  \draw[otMuted!60,line width=0.4pt] (0,0) grid (4,4);
  \draw[otMuted,line width=0.7pt,dash pattern=on 2.5pt off 1.8pt] (0,0) -- (4.35,4.35);
  \draw[otGreen!60!black,line width=0.7pt,dash pattern=on 2.5pt off 1.8pt] (-0.45,0.55) -- (3.45,4.45);
  \draw[otNavy,line width=1.6pt,line join=round] (0,0) -- (1,0) -- (1,1) -- (1,2) -- (2,2) -- (2,3);
  \draw[otGreenDark,line width=1.3pt,dash pattern=on 3pt off 1.8pt,line join=round]
    (2,3) -- (3,3) -- (3,4) -- (4,4);
  \fill[otRed!85!black] (1,2) circle (2.6pt);
  \node[font=\scriptsize,text=otNavy,anchor=east] at (-0.15,-0.3) {$(0,0)$};
  \node[font=\scriptsize,text=otMuted,anchor=west] at (4.42,4.28) {$y=x$};
  \node[font=\scriptsize,text=otGreen!50!black,anchor=south] at (2.8,4.5) {$y=x+1$};
\end{tikzpicture}\\[0.5mm]
\legende{Un chemin fautif touche $y=x+1$ (point rouge) : la réflexion de
son début l'identifie à un chemin ordinaire — c'est la bijection de
comptage.}
\end{center}
\end{corrige}
\end{blocLivre}

\begin{blocLivre}{cor-16-22}
\begin{corrige}{22}{Identités par double comptage}
\demarche{Le même ensemble, deux inventaires : l'égalité est gratuite.}
\begin{questions}
  \item Coupons $2n$ personnes en deux moitiés de $n$. Une paire est
    interne à la première moitié ($\binom n2$ choix), interne à la
    seconde ($\binom n2$), ou mixte ($n\times n$) :
    $\res{\dbinom{2n}2=2\dbinom n2+n^{2}}$ — ce qui se vérifie aussi par
    le calcul.
  \item Classons les parties $\ens{a<b<c}$ de $\ens{1,\dots,n+1}$ selon
    leur maximum $c=k+1$ : il reste à choisir $\ens{a,b}$ parmi les $k$
    entiers inférieurs, soit $\binom k2$ choix. En sommant :
    $\res{\displaystyle\sum_{k=2}^{n}\binom k2=\binom{n+1}{3}}$ — la
    « crosse de hockey » du triangle de Pascal.
\end{questions}
\end{corrige}
\end{blocLivre}

\begin{blocLivre}{cor-16-23}
\begin{corrige}{23}{Permutations et points fixes}
\demarche{Choisir les points fixes, déranger le reste — puis compter les
couples (permutation, point fixe).}
\begin{questions}
  \item Choisir les $k$ points fixes ($\binom nk$ façons) puis déranger
    les $n-k$ autres ($D_{n-k}$ façons) :
    $\res{\dbinom nk D_{n-k}}$ permutations à exactement $k$ points
    fixes.
  \item Ces classes partitionnent $S_n$ :
    $\sum_k\binom nk D_{n-k}=n!$. Comptons maintenant les couples
    $(\sigma,i)$ avec $\sigma(i)=i$ : à $i$ fixé, $(n-1)!$ permutations
    conviennent, d'où $n\times(n-1)!=n!$ couples ; le nombre moyen de
    points fixes vaut donc $\dfrac{n!}{n!}=\res{1}$ — quelle que soit la
    taille $n$ : un avant-goût de l'espérance (chapitre 17).
\end{questions}
\end{corrige}
\end{blocLivre}

\begin{blocLivre}{cor-16-24}
\begin{corrige}{24}{Fibonacci compte aussi}
\demarche{Le dernier élément entre ou n'entre pas : Fibonacci apparaît ;
un décalage compte ensuite chaque taille.}
\begin{questions}
  \item $a_1=2$, $a_2=3$, $a_3=\res{5}$. Si $n$ est dans la partie, $n-1$
    n'y est pas et le reste est une partie admissible de
    $\ens{1,\dots,n-2}$ ; sinon c'en est une de $\ens{1,\dots,n-1}$ :
    $a_n=a_{n-1}+a_{n-2}$. Avec $a_1=2=F_3$ et $a_2=3=F_4$ :
    $\res{a_n=F_{n+2}}$ (chapitres 5, 12 et 13).
  \item Pour une partie $\ens{x_1<\dots<x_k}$ sans éléments consécutifs,
    poser $y_i=x_i-(i-1)$ : les $y_i$ sont strictement croissants dans
    $\ens{1,\dots,n-k+1}$, sans autre contrainte — bijection avec les
    parties à $k$ éléments : $\dbinom{n-k+1}{k}$ choix. En sommant :
    \[
    \res{\sum_{k\geqslant0}\binom{n-k+1}{k}=F_{n+2}} :
    \]
    les diagonales du triangle de Pascal se somment en nombres de
    Fibonacci.
\end{questions}
\end{corrige}
\end{blocLivre}

% !TEX root = ../../main.tex
%% Corrigé du défi du chapitre 16 — vérifié par verifications/defis_16_20.py.

\begin{blocLivre}{cor-16-25a}
\begin{corrige}{25}{Les nombres de Stirling de seconde espèce}
\demarche{Tout repose sur deux idées : suivre le sort d'un élément
particulier (récurrences), et classer des applications selon leur image
(identités).}
\partieExo{A}{Premières valeurs et récurrence}
\begin{questions}
  \item $\res{S(n,1)=S(n,n)=1}$ (un seul bloc ; que des singletons). Avec
    $n-1$ blocs, un seul bloc a deux éléments, les autres sont des
    singletons : $\binom n2$ choix de la paire. Avec deux blocs $\ens{A,\overline A}$ :
    les couples $(A,\overline A)$ avec $A\neq\varnothing,\intFF1n$ sont au
    nombre de $2^{n}-2$, chaque partition étant comptée deux fois :
    $\res{S(n,2)=2^{n-1}-1}$.
  \item Si $\ens{n}$ est un bloc, le reste est une partition de
    $\intFF1{n-1}$ en $k-1$ blocs : $S(n-1,k-1)$ cas. Sinon, $n$ rejoint
    l'un des $k$ blocs d'une partition de $\intFF1{n-1}$ en $k$ blocs :
    $k\,S(n-1,k)$ cas. D'où la relation. Ligne $5$ :
    $1,\ 1+2\cdot7=15,\ 7+3\cdot6=25,\ 6+4\cdot1=10,\ 1$ ; ligne $6$ :
    $1,\ 31,\ 90,\ 65,\ 15,\ 1$.
  \item $S(4,2)=2^{3}-1=\res{7}$ partitions (figure) : trois de type
    « $2+2$ » et quatre de type « $1+3$ ».
\end{questions}
\begin{center}
\begin{tikzpicture}[line cap=round]
  \foreach \blocA/\blocB [count=\i] in {{1}/{2,3,4},{2}/{1,3,4},{3}/{1,2,4},{4}/{1,2,3},{1,2}/{3,4},{1,3}/{2,4},{1,4}/{2,3}}{%
    \begin{scope}[xshift={(\i-1)*1.85cm}]
      \coordinate (p1) at (0,0.8); \coordinate (p2) at (0.8,0.8);
      \coordinate (p3) at (0.8,0); \coordinate (p4) at (0,0);
      \foreach \a in \blocA {\foreach \b in \blocA {%
        \ifnum\a<\b \draw[otGreenDark,line width=1.4pt] (p\a) -- (p\b);\fi}}
      \foreach \a in \blocB {\foreach \b in \blocB {%
        \ifnum\a<\b \draw[otNavy,line width=1.4pt] (p\a) -- (p\b);\fi}}
      \foreach \a in \blocA {\filldraw[fill=otGreen,draw=otGreenDark] (p\a) circle (2.4pt);}
      \foreach \a in \blocB {\filldraw[fill=otNavy!35,draw=otNavy] (p\a) circle (2.4pt);}
      \foreach \k in {1,2,3,4}{\node[font=\tiny,text=otMuted] at ($(p\k)+({\k==2||\k==3 ? 0.22 : -0.22},{\k<3 ? 0.18 : -0.18})$) {\k};}
    \end{scope}}
\end{tikzpicture}\\[0.5mm]
\legende{Les sept partitions de $\ens{1,2,3,4}$ en deux blocs : les points
reliés sont dans le même bloc (un bloc en vert, l'autre en bleu).}
\end{center}
\end{corrige}
\end{blocLivre}

\begin{blocLivre}{cor-16-25b}
\begin{corrigeSuite}
\partieExo{B}{Surjections et factorielles descendantes}
\begin{questions}[start=4]
  \item Les fibres d'une surjection $f:\intFF1n\to\intFF1k$ forment une
    partition en $k$ blocs, étiquetés par les $k$ valeurs distinctes ;
    réciproquement, une partition et une façon d'étiqueter ses blocs
    ($k!$ façons) définissent une surjection : $\res{k!\,S(n,k)}$
    surjections.
  \item Il y a $x^{n}$ applications de $\intFF1n$ dans $\intFF1x$. Celles
    dont l'image a $j$ éléments s'obtiennent en choisissant l'image
    ($\binom xj$ façons) puis une surjection sur elle ($j!\,S(n,j)$) ; et
    $\binom xj\,j!=x(x-1)\cdots(x-j+1)$, nul si $j>x$ : on obtient
    l'identité.
  \item Les deux membres sont des polynômes en $x$ de degré au plus $n$, égaux
    en tout entier naturel, donc en une infinité de points : ils sont
    égaux. En divisant par $j!$ : $\res{X^{n}=\sum_jj!\,S(n,j)\,H_j}$. Pour
    $n=3$ : $1!\cdot1$, $2!\cdot3$, $3!\cdot1$, soit
    $X^{3}=H_1+6H_2+6H_3$ — exactement la diagonale de la table des
    différences du chapitre 12.
\end{questions}
\partieExo{C}{Nombres de Bell et formule explicite}
\begin{questions}[start=7]
  \item Le bloc de $n+1$ s'écrit $\ens{n+1}\cup T$, $T\subseteq\intFF1n$ ;
    si les $k$ éléments restants sont partagés librement ($B_k$ façons),
    $T$ a $n-k$ éléments, choisis de $\binom n{n-k}=\binom nk$ façons :
    $\res{B_{n+1}=\sum_{k=0}^{n}\binom nkB_k}$. On trouve
    $B_1,\dots,B_6=1,2,5,15,52,\res{203}$ — sommes des lignes du triangle.
  \item D'après l'exercice 15, il y a $3^{n}-3\cdot2^{n}+3$ surjections sur
    trois éléments, et c'est $3!\,S(n,3)$ (question 4) :
    $\res{S(n,3)=\dfrac{3^{n}-3\cdot2^{n}+3}{6}}$. Contrôle :
    $n=4$ donne $\frac{81-48+3}{6}=6$, $n=5$ donne $\frac{243-96+3}{6}=25$.
\end{questions}
\end{corrigeSuite}
\end{blocLivre}

